\chapter{Skills Acquired}
\label{ch:skills}

This chapter catalogues the skills I acquired during my internship at
FlowGenX AI. They fall into two broad categories: technical skills,
which I developed through the day-to-day work on the
\texttt{runtime\_connectivity} project and the surrounding repositories,
and non-technical skills, which I developed through working inside a
remote-first, globally distributed engineering team.

\section{Technical Skills}

\subsection{Programming Languages}
\textbf{Python 3.11+} was the primary language of the internship and
the one I now use most fluently. Across five and a half months and
178 commits I wrote idiomatic, typed, async-aware Python: Pydantic
models, decorators, context managers, generators, dataclasses,
asyncio coroutines, pytest fixtures, and mocking patterns. Black,
isort, Ruff, and pre-commit hooks made the ``correct shape'' of
clean Python second nature. \textbf{SQL} got daily exposure across
PostgreSQL, MySQL, SQL Server, and the cloud managed-database
variants (AWS RDS, Azure SQL): CRUD, schema inspection, bulk-write
logic, and the parameter standardisation work that unified how SQL
connectors accept both DataFrames and raw lists of records.
\textbf{TypeScript / JavaScript} stayed at read-level: while my own
contributions were primarily backend, I regularly read the Next.js
frontend to understand how the UI consumed the auth-config and
OpenAPI metadata my connectors produced. \textbf{Configuration
languages} (JSON, YAML, TOML) were used pervasively for OpenAPI
specs, Docker Compose, Kubernetes manifests, and Python project
files.

\subsection{Frameworks and Libraries}
\textbf{FastAPI} and \textbf{Pydantic} sit at the centre of the
platform's backend; I worked with FastAPI patterns through the
runtime APIs and used Pydantic extensively to model connection
parameters and request/response schemas for every connector.
\textbf{Apache Airflow} is the workflow engine underneath
\texttt{runtime\_apis}, and I learned how DAGs are defined,
scheduled, and triggered, and how the connectors I shipped become
callable nodes inside them. \textbf{Pandas} drove the in-memory
DataFrame transformations in the database and file connectors, and
the DataFrame/list duality I introduced into the SQL write methods.
\textbf{boto3} was used pervasively across the AWS connectors (S3,
RDS, Textract, Rekognition) including paginators, signed-request
flows, multipart uploads, asynchronous job patterns, and IAM
credential handling. Other libraries used routinely include
\texttt{requests}, \texttt{httpx}, \texttt{asyncio},
\texttt{pytest} (with fixtures, mocks, and integration markers),
and \texttt{Jinja2}.

\subsection{Databases}
Hands-on familiarity with a wide range of database backends through
the connectors I built and tested:

\begin{itemize}[leftmargin=*]
    \item \textbf{Relational:} PostgreSQL, MySQL, Microsoft SQL
    Server, Azure SQL Server, AWS RDS PostgreSQL, AWS RDS MySQL,
    Snowflake.
    \item \textbf{Document and NoSQL:} MongoDB.
    \item \textbf{Search:} Elasticsearch, Typesense.
    \item \textbf{Vector databases:} Pinecone, Weaviate, Milvus.
\end{itemize}

For each I learned the connection model, the authentication options,
the standard CRUD shape, the pagination model, and the
performance-relevant operations (bulk insert, upsert, truncate).

\subsection{Cloud Platforms}
On \textbf{AWS} I worked directly with S3 (bucket policies,
versioning, multipart uploads), RDS for both PostgreSQL and MySQL,
Textract (sync and async document analysis with adapters),
Rekognition (image, video, face, custom labels, stream processors),
and IAM credential flows. On \textbf{Azure} I worked with Blob
Storage, Data Lake, and Azure SQL Server (with firewall-rule mode
selection and retry logic). At the connector level I also built
work on \textbf{Google Cloud} services (Tasks, Forms, Chat, Meet)
and \textbf{Microsoft~365} (OneDrive and SharePoint).

\subsection{AI / LLM Providers and Generative-AI Patterns}
I built and tested provider connectors for \textbf{OpenAI}
(assistants, models, threads), \textbf{Anthropic} (messages, batch,
files, models, the Skills API), \textbf{Google Gemini},
\textbf{xAI}, \textbf{Groq}, \textbf{DeepSeek}, and
\textbf{Perplexity AI}. The \textbf{vector-database} connectors
(Pinecone, Weaviate, Milvus) gave me working familiarity with
embeddings, similarity search, metadata filtering, and the role
these systems play inside RAG pipelines. Working inside FlowGenX
also exposed me to the practical use of multi-agent patterns
(ReAct, Supervisor, Swarm, Deep) and the underlying protocols ---
the \emph{Model Context Protocol (MCP)} for exposing tools to
agents and \emph{Agent-to-Agent (A2A)} communication for inter-agent
collaboration.

\subsection{Authentication Patterns}
The connector work exposed me to every modern authentication pattern
in production use:

\begin{itemize}[leftmargin=*]
    \item \textbf{OAuth 2.0} with refresh-token flows (the most
    important being the QuickBooks access-token refresh mechanism I
    implemented).
    \item \textbf{API Key} and \textbf{Bearer Token} authentication.
    \item \textbf{HTTP Basic Authentication}.
    \item \textbf{AWS IAM} credentials and signed-request flows.
\end{itemize}

\subsection{System Architecture and Concepts}
FlowGenX is composed of multiple cooperating services
(\texttt{flowgenx-nextjs-V2}, \texttt{runtime\_apis},
\texttt{runtime\_components}, \texttt{runtime\_connectivity}), and
working inside this \textbf{microservice architecture} gave me a
concrete intuition for service boundaries, package-level
interfaces, and the discipline required to evolve a shared library
without breaking its consumers. Every connector method I wrote
produced \textbf{OpenAPI 3.0} metadata via the
\texttt{@openapi\_endpoint} decorator, so by the end of the
internship I could read and write OpenAPI schemas fluently. I also
implemented \textbf{webhook} configuration for connectors such as
Apify and participated in the broader event-driven model of the
runtime, including the \emph{Webhook Trigger} that initiates
workflow execution when an external system posts data. \textbf{gRPC}
stayed at read-level through the runtime APIs codebase, and I
worked with \textbf{Docker} for local development, observed the
\textbf{Kubernetes / Helm / Terraform} setup the platform uses in
production, and was directly exposed to caching patterns through
the connector cleanup work and the SQL Server retry logic.

Beyond the connector layer, I gained a working understanding of
several \textbf{platform-level products} that surround it:
\emph{Agent Fabric} (multi-agent orchestration runtime),
\emph{Agent Gateway} (governance/security plane), \emph{VectorVerse}
(vector-powered Knowledge Base), the \emph{Playground} (interactive
testing surface for MCP, agent, and workflow modes), the \emph{API
Workspace} (unified management for OpenAPI tools and
authentication), and the \emph{Data Compiler} (schema-drift
component that keeps long-running pipelines resilient to upstream
change) --- giving me an end-to-end mental model of how an agentic
platform fits together.

\subsection{Testing}
I wrote hundreds of \textbf{pytest} cases across the connectors I
delivered (fixtures, mocks, parameterised tests, integration
markers), plus a dedicated \textbf{live-API integration test
script} for every connector that produced a
\textsc{PASS}/\textsc{FAIL}/\textsc{SKIP} report against the real
third-party API. Every connector was finally validated through
the platform's API tester, MCP Playground, and end-to-end workflow
execution.

\subsection{Tooling and IDEs}
\begin{itemize}[leftmargin=*]
    \item \textbf{Git and GitHub} --- feature-branch workflow, pull
    requests, code review, merge conflict resolution.
    \item \textbf{ClickUp} --- task tracking, status reporting, and
    sprint visibility.
    \item \textbf{uv} --- the modern, fast Python package manager that
    FlowGenX standardises on.
    \item \textbf{Black, isort, Ruff, pre-commit} --- the code-quality
    toolchain that enforces uniform style across the codebase.
    \item \textbf{Docker and Docker Desktop} --- local container
    development.
    \item \textbf{Visual Studio Code} --- my primary IDE for the
    internship, with Python, Pylance, and Docker extensions.
    \item \textbf{Postman} --- exploration of third-party APIs before
    coding the connector.
    \item \textbf{Google Chat and Google Meet} --- the daily
    communication tools.
\end{itemize}

\section{Non-Technical and Professional Skills}

\subsection{Communication and Code Review}
Working as the only intern on a remote, distributed team taught me to
write more carefully: self-contained status updates, asynchronous
phrasing of technical questions, blockers raised early in writing
rather than in synchronous meetings, and questions asked in public
channels (so answers benefited the team and stayed searchable). On
\textbf{code review}, every pull request I opened received review
comments, and I came to appreciate review as one of the
highest-value activities a team can do --- treating each comment as
a conversation rather than a verdict, pushing back politely with
reasoning when I disagreed, and thanking reviewers for catching
subtle issues. By the second half of the internship I was
occasionally reviewing teammates' pull requests myself, a small but
meaningful sign of having earned the team's trust.

\subsection{Time Management, Self-Direction, and Cross-Time-Zone Work}
A remote internship with flexible hours requires more self-discipline,
not less. I built the habit of setting a daily intention --- ``ship
one connector to a reviewable state, fix one bug, or close one
documentation deliverable'' --- and stayed honest with myself about
what counted as ``done''. The ClickUp board and daily pull-request
flow made this discipline tangible. Working across the
roughly twelve-hour Dhaka-to-San~Francisco offset taught me to leave
detailed end-of-day updates so the United States team could pick up
where I left off, to compose questions that did not require immediate
replies, and to schedule synchronous meetings during the evening
overlap window without disrupting either side's focus time.

\subsection{Professionalism, Curiosity, and Resilience}
The defining cultural expectation at FlowGenX --- as captured in my
offer letter --- was that interns demonstrate the same commitment and
work ethic as full-time team members, and I tried to honour that
throughout: every commit followed the convention, every pull request
was self-reviewed before opening, every connector was tested before
being declared complete, every deliverable documented for the next
person who would touch it. The breadth of technologies on display
(90+ third-party APIs, multiple cloud providers, several agentic
patterns, MCP, A2A, Airflow) forced a kind of structured curiosity ---
reading API documentation efficiently, skimming a codebase for
patterns rather than reading every line. And every connector exposed
at least one quirky issue (a malformed base URL, an undocumented
rate limit, an auth flow that worked on Linux but not Windows),
which I learned to treat as a puzzle with a finite solution rather
than a blocker, and to document so the next person hitting the same
issue would not start from scratch. The OpenAPI sample-backfill work
was the clearest example of patience and attention to detail:
replacing placeholder values with real ones across hundreds of
endpoints is unglamorous but disproportionately impactful for the
end-user experience.
